\documentclass[a4paper,twoside]{article}

\usepackage[LGR,T1]{fontenc}
\usepackage[utf8x]{inputenc}
\usepackage[english]{babel}
\usepackage{graphicx}
\usepackage{lipsum}  

\usepackage{epsfig}
\usepackage{subcaption}
\usepackage{calc}
\usepackage{amssymb}
\usepackage{amstext}
\usepackage{amsmath}
\usepackage{amsthm}
\usepackage{multicol}
\usepackage{pslatex}
\usepackage{natbib}
\usepackage{apalike}
\usepackage{xurl}
\usepackage[algoruled,lined,linesnumbered,algosection]{algorithm2e}
\usepackage{booktabs}
\usepackage{dsfont}
\usepackage{pifont} 
\usepackage{SCITEPRESS}


\newtheorem{defn}{Definition}

\begin{document}

\title{Scored Rule Sets for Interpretable Multiclass Classification}


%\author{\authorname{Robin Nunkesser\orcidAuthor{0009-0000-1238-8987}} \affiliation{Hamm-Lippstadt University of Applied Sciences, Marker Allee 76--78, 59063 Hamm, Germany}\email{robin.nunkesser@hshl.de}}

\author{\authorname{Anonymous} \affiliation{Affiliation}\email{anonymous@example.com}}

\keywords{Machine Learning, Classification, Interpretability, Genetic Programming, Research Siloing, Hypothesis Space}

\abstract{Interpretable multiclass classification is studied across several research communities, but comparisons and conceptual integrations across rule sets, rule lists, decision trees, and evolutionary approaches are still limited. This paper addresses this gap by proposing scored rule sets as a unifying and interpretable hypothesis space. We first analyze structural relationships among established algorithms and identify indications of research siloing, especially between evolutionary and non-evolutionary lines of work. We then compare selected highly interpretable multiclass learners on UCI benchmark datasets. Beyond empirical comparison, we provide a formal definition of scored rule sets, show how they generalize classical rule-based models, and illustrate concrete transformations from logicGP, ExSTraCS, and CART representations. The resulting perspective clarifies when scored rule sets preserve interpretability while increasing modeling flexibility. Overall, the results support scored rule sets as a practically useful bridge between existing interpretable model families and as a basis for future cross-community method development.}



\onecolumn \maketitle \normalsize \setcounter{footnote}{0} \vfill

\section{\uppercase{Introduction}}

Classification is a fundamental task in machine learning, where the goal is to assign labels to instances based on their features. The choice of classification algorithm can significantly impact the performance and interpretability of the resulting model. The landscape of classification algorithms is vast, encompassing a wide range of approaches from simple rule-based or linear models to complex ensemble methods and deep learning architectures. 

%In this paper, we focus on highly interpretable classification algorithms, with a particular emphasis on multiclass classification, which are common in many real-world applications.

According to \citet{RudinEtAlSurvey2022}, interpretability is a crucial aspect of machine learning models, especially in high-stakes decision-making scenarios. Interpretable models allow practitioners to understand the reasoning behind predictions, which can be essential for building trust and ensuring ethical use of machine learning. In the context of multiclass classification, interpretability becomes even more important, as the decision boundaries can be more complex and harder to understand.

For tabular data, which is prevalent in many applications, rule- or decision-based models are often considered highly interpretable. In addition, score-based models that provide insights into feature importance can also be valuable for interpretability \citep{RudinEtAlSurvey2022}. 

Interpretable classification algorithms have been developed in various research communities, including machine learning, data mining, bioinformatics, and evolutionary computation. However, there is a lack of comprehensive comparisons of these algorithms. An analysis of the relationships between these algorithms suggests that there may be a certain amount of research siloing, with limited integration of knowledge across different communities, especially in the case of rule-set-based algorithms and those that utilize evolutionary computation techniques.

This paper aims to fill this gap by providing a unified hypothesis space for many highly interpretable classification algorithms and conducting a comparative analysis of highly interpretable classification algorithms for multiclass problems, with a particular focus on those that utilize evolutionary computation techniques. The main contributions of this paper are:
\begin{itemize}
  \item A unifying formalization of scored rule sets as an interpretable hypothesis space that generalizes classical rule sets and can encode rule-list and decision-tree structures.
  \item A cross-community analysis of algorithm relationships that highlights potential research siloing between evolutionary and non-evolutionary interpretable classification approaches.
  \item An empirical comparison on selected UCI multiclass benchmarks, combined with concrete model transformations (logicGP, ExSTraCS, CART) that demonstrate practical applicability of the unifying representation.
\end{itemize}

The remainder of this paper is organized as follows: we first review interpretable hypothesis spaces and their relations, then analyze potential cross-community siloing, and afterwards present the experimental comparison. We then discuss the empirical findings and, building on these observations, formalize scored rule sets and their relation to existing model classes. This discussion-before-formalization order is intentional: we first motivate the unifying view from empirical observations and then introduce the formal model-space definition.

The practical importance of this unifying view is that it enables method-agnostic comparison and interpretation across model families that are usually evaluated in isolation. In particular, it provides a common language for model structure, prediction aggregation, and compactness-oriented analysis.


\section{\uppercase{Interpretable Hypothesis Spaces}}

The hypothesis space of an algorithm defines the set of all possible models that can be learned by the algorithm. The choice of hypothesis space is crucial for the performance and interpretability of the resulting model. 

\citet{RudinEtAlSurvey2022} provide a comprehensive survey of interpretable machine learning methods. For tabular data, which is prevalent in many applications, rule- or decision-based models are often considered highly interpretable. In addition, score-based models that provide insights into feature importance can also be valuable for interpretability. %As many modern rule- or decision-based models include score-based components, we will focus on these models in our analysis (which are a generalization of pure rule- or decision-based models).

This analysis focuses on single interpretable models and largely excludes approaches based on extensive ensembling or stacking, as such methods generally diminish model transparency. Importantly, \citet{RudinEtAlSurvey2022} demonstrate that this restriction need not entail a loss of predictive performance. With respect to decision trees, they explicitly note:

\begin{quote}
  ``When fully optimized, single trees can be as accurate as ensembles of trees, or neural networks, for many problems.''
\end{quote}

\citet{RudinEtAlSurvey2022} list three rule- or decision-based models that are designed for interpretability and can be used for classification: rule sets, rule lists, and decision trees. Figure~\ref{fig:hypospaces} shows examples of binary classification problems for highly interpretable hypothesis spaces \citep[adapted from][]{imodels2021}. The first column shows an example of a rule set, the second column shows an example of a rule list, and the third column shows an example of a rule / decision tree.

\begin{figure*}
  \centering
  \includegraphics[width=\textwidth]{contours}
  \caption{Binary classification examples for highly interpretable hypothesis spaces}
  \label{fig:hypospaces}
\end{figure*}


% \begin{aligned}
% \textbf{IF } X_1 < 6 \text{ and } X_2 > 6: \quad & p \mathrel{+}= 0.8 \\
% \textbf{IF } X_1 < 5 \text{ and } X_2 < 7: \quad & p \mathrel{+}= 0.3 \\
% \textbf{IF } X_1 > 4 \text{ and } X_2 < 4: \quad & p \mathrel{+}= 0.8
% \end{aligned}

% \begin{aligned}
% \textbf{IF } X_1 < 5: \quad & p = 0.3 \\
% \textbf{ELSE IF } X_2 < 4: \quad & p = 0.9 \\
% \textbf{ELSE IF } X_1 > 6: \quad & p = 0.7 \\
% \textbf{ELSE}: \quad & p = 0.5
% \end{aligned}



%\subsection{Binary Classification}

\subsection{Overview of Rule-Based Models}

In the following, we will give an overview of the different types of highly interpretable classification algorithms based on rules or decisions, which are commonly used for tabular data. We will focus on rule sets, rule lists, and decision trees, as these are the most common types of rule-based models. We will also discuss the advantages and disadvantages of each type of model and their applicability to multiclass classification problems.

\paragraph{Rule Sets}

Rule sets are a collection of rules that can be applied in any order to classify instances. Each rule consists of a condition and a corresponding score or class label. Rule sets can be highly interpretable, as each rule can be easily understood and analyzed by humans. They are closely related to disjunctive normal form (DNF) representations of Boolean functions, which can represent any Boolean function as a disjunction of conjunctions. The $A^q$ algorithm \citep{1571135649817801472} is a good example of an algorithm originating in logic synthesis that was adapted for classification in machine learning \citep{michalski1973aqval}.  

Modern examples of rule set algorithms from different communities include RIPPER \citep{10.5555/3091622.3091637}, Slipper \citep{10.5555/315149.315320}, M5Rules \citep{10.1007/3-540-46695-9_1}, GPAS \citep{10.1093/bioinformatics/btm522}, RuleFit \citep{friedman_predictive_2008}, MLRules \citep{10.1145/1390156.1390185}, BioHEL \citep{bacardit_improving_2009}, ExSTraCS \citep{urbanowicz_exstracs_2015}, RuleKit \citep{gudys_rulekit_2020}, and logicGP \citep{10.1145/3712255.3734300}.

\paragraph{Rule Lists}

Rule lists are an ordered sequence of rules, where each rule is applied in a specific order to classify instances. Each rule consists of a condition and a corresponding score or class label. Rule lists can be highly interpretable, as the order of the rules can provide insights into the decision-making process. \citet{QUINLAN1987221} states that they have the potential to simplify overly complex decision trees.

Modern examples of rule list algorithms from different communities include  OneR \citep{holte_very_1993} and BRL \citep{letham_interpretable_2015}.

\paragraph{Decision Trees}

Decision trees are a popular choice for interpretable classification models. They represent the decision-making process as a tree structure, where each internal node represents a decision based on a feature, and each leaf node represents a class label. Decision trees can be easily visualized and understood by humans, making them highly interpretable. 

Modern examples of decision tree algorithms from different communities include
optimized versions of CART \citep{Breiman1984} and C4.5 \citep{quinlan1993}, DL8.5 \citep{Aglin_Nijssen_Schaus_2020}, GOSDT \citep{lin2022generalizedscalableoptimalsparse}, and logicDT \citep{lau_logicdt_2024}. Notable post-hoc methods for optimizing decision trees include TAO \citep{NEURIPS2018_185c29dc} and HS \citep{pmlr-v162-agarwal22b}.

\subsection{Multiclass Classification}

A very general approach to multiclass classification is to use strategies like One-Versus-All or One-Versus-One, which can be applied to any binary classification algorithm. However, these strategies can lead to less interpretable models, as they require multiple binary classifiers to be trained and combined to make a final prediction.

As an alternative, some algorithms are designed to directly model multiclass classification problems without relying on such strategies. These algorithms can provide more interpretable models, as they can capture the relationships between classes and features in a more direct way. For rule lists and decision trees, this can be achieved by directly allowing different classes in the leaf nodes / rules. For rule sets, this cannot be achieved directly because rules can be applied in any order, which can lead to ambiguity in the decision-making process. However, this can be achieved by either using a voting scheme, where each rule votes for a class and the class with the most votes is selected as the final prediction, or by using a scoring scheme, where each rule is associated with a score for each class and the class with the highest cumulative score is selected as the final prediction. This can also enhance interpretability, as the scores can provide insights into the importance of each rule in the decision-making process. This approach can also be used for rule lists and decision trees, where each rule or leaf node can be associated with a score for each class. This allows for a more flexible and interpretable modeling of multiclass classification problems, as it can capture the relationships between classes and features in a more direct way.

\subsection{Comparison}

Rule lists and decision trees are very similar in terms of their structure and interpretability. Rule lists can be seen as a special case of decision trees, where the tree is linear. Rule sets, on the other hand, are more flexible and can represent a wider range of decision boundaries, which is a very desirable property for modeling complex relationships in data. It is easy to see that rule sets can represent any decision tree or rule list \citep[see e.g.][]{witten2025data}. It is also possible to construct a decision tree from a rule set, however the replicated subtree problem may occur \citep{witten2025data}. \citet{fuernkranz2012foundations} list additional advantages of rule sets compared to rule lists and decision trees.

Potential disadvantages of rule sets are:

\begin{itemize}
  \item contradictory rules can occur in rule sets, which can lead to ambiguity in the decision-making process, 
  \item some instances may not be covered by any rule, which can lead to unclassified instances, and
  \item they can be less interpretable than rule lists and decision trees, as the rules can be applied in any order, which can make it harder to understand the decision-making process,
  \item they can be more computationally expensive to learn than decision trees, as the search space for rule sets is larger than that for decision trees.  
\end{itemize}

We will address these disadvantages later by using scored rule sets, which can provide a more flexible and interpretable modeling of multiclass classification problems. Before that, we will discuss the potential siloing of research communities in the field of rule-based classification algorithms and the need for knowledge integration across different communities. The focus will be on the intersection of rule-based classification algorithms and evolutionary computation techniques, as this is an area where there seems to be a certain amount of research siloing.

\section{\uppercase{Potential Siloing of Research Communities}}

Rule- and decision-based classification is studied across a broad range of disciplines within Statistics and Computer Science, including machine learning, data mining, bioinformatics, and evolutionary computation. Research in this area is distributed across largely separate communities, and no single comprehensive survey covers all existing approaches. Within the evolutionary computation literature, the review by \citet{urbanowicz_moore_2009} provides a thorough treatment of rule-based classification methods that employ evolutionary techniques, yet it does not encompass approaches from other communities. \citet{alcala-fdez_keel_2009} present a cross-community comparison via the KEEL framework; however, development of that framework ceased around 2015, and more recent algorithms are therefore not represented.

%\subsection{Related Work and Positioning}

Related work can be grouped into three streams. First, broad interpretable-ML surveys such as \citet{RudinEtAlSurvey2022} discuss model families and design principles, but do not provide a focused cross-community synthesis of multiclass rule-based approaches. Second, domain-specific overviews in evolutionary computation, e.g., \citet{urbanowicz_moore_2009}, provide substantial depth for LCS-style methods but are naturally centered on that community. Third, framework-based comparisons such as \citet{alcala-fdez_keel_2009} improve comparability across methods, yet are constrained by the algorithm sets available in their tooling ecosystems.

This paper is positioned at the intersection of these streams: we combine cross-community structuring, an empirical comparison on shared benchmarks, and a unifying scored-rule-set formalization that directly connects rule sets, rule lists, and decision trees. In this sense, our contribution is not a new learner, but a unifying model-space perspective with concrete transformations and comparative evidence.

Compared with prior work, the central added value is this explicit cross-family mapping into one interpretable hypothesis space, which makes methodological differences analyzable without changing the target representation.

A literature review of the algorithms considered in this paper reveals evidence of research siloing and indicates that cross-community knowledge integration would be beneficial (see Figure~\ref{fig:sjobergsilos}, adapted from \cite{10.1007/978-3-319-49094-6_1,software5010012}, for an illustration of knowledge hierarchy under siloed research conditions).

\begin{figure}
  \centering
  \includegraphics[width=\columnwidth]{sjobergsilos2}  
  \caption{Knowledge hierarchy within the context of siloed research (adapted from \cite{software5010012})}
  \label{fig:sjobergsilos}
\end{figure}

The present work focuses specifically on the intersection of rule-based classification and evolutionary computation. Figure~\ref{fig:hicitations} depicts the citation relationships among highly interpretable classification algorithms, with particular emphasis on those employing evolutionary techniques. Algorithms are grouped by their underlying methodology; edges denote the strength of the relationship, with solid lines indicating strong relationships (extension or direct inspiration) and dashed lines indicating medium relationships (comparison).

\begin{figure*}
  \centering
  \includegraphics[width=.65\textwidth]{hicitations}  
  \caption{Relationships between highly interpretable classification algorithms}
  \label{fig:hicitations}
\end{figure*}

Figure~\ref{fig:hicitations} shows that evolutionary computation techniques have contributed to the development of several highly interpretable classification algorithms; however, cross-community connectivity remains limited. Two structural observations stand out:

\begin{itemize}
  \item No citation edges lead from non-evolutionary algorithms to evolutionary algorithms, suggesting that the latter are largely unacknowledged by the former community.
  \item Where connections from evolutionary algorithms to non-evolutionary ones exist, they predominantly reference foundational works; propagation through chains of extensions and inspirations has not resulted in sustained integration between the two communities.
\end{itemize}

These observations indicate that meaningful opportunities exist for cross-methodological integration in interpretable classification research. Furthermore, evaluations of several evolutionary algorithms included comparisons with methods that are neither designed for interpretability nor rule-set-based. For instance, logicGP \citep{10.1145/3712255.3734300} was evaluated against the full set of ML.NET AutoML classification algorithms.

\section{\uppercase{Experiments}}

The result of the potential siloing of research communities is that there is a lack of comprehensive comparisons of highly interpretable classification algorithms, especially those that utilize evolutionary computation techniques.  
To address this gap, we will conduct a series of experiments. The goal is to assess their performance in terms of F1 score, interpretability, and computational efficiency.

Candidate algorithms for comparison have to fulfill the following criteria: 

\begin{itemize}
  \item they have to be able to directly model multiclass classification problems, %without relying on strategies like One-Versus-All or One-Versus-One, which can lead to less interpretable models,
  \item they have to be available as open source, allowing for a fair and transparent comparison, and
  %\item they have to be applicable to tabular data, which is the focus of our analysis,
  \item they have to be designed for interpretability, meaning that the resulting models should be easily understandable by humans, 
  %\item they have to be competitive in terms of performance.
\end{itemize}

Table~\ref{tab:algos} gives an overview of some algorithms mentioned in the previous sections that are candidates for comparison in our experiments. The table includes information about the year of publication, whether the algorithm can directly model multiclass classification problems, and whether it has open source implementations available in popular machine learning libraries such as scikit-learn, R, Weka, and ML.NET. The algorithms are categorized into four groups: Rule Tree, Rule List, Rule Set, and Evolutionary Computation.

\begin{table}
  \centering
  \caption{Overview of Algorithms (MC: Multiclass, sl: scikit-learn, R: R, W: Weka, MN: ML.NET)}
  \label{tab:algos}
  \begin{tabular}{p{1.5cm}p{0.6cm}p{0.33cm}p{0.33cm}p{0.33cm}p{0.33cm}p{0.33cm}}
    \toprule
    Algorithm & Year & MC & sl & R & W & MN \\
    \midrule
    \multicolumn{7}{l}{\textit{Rule Tree}} \\
    \midrule
    CART & 1984 & \checkmark & \checkmark & \checkmark & \checkmark & \ding{55} \\
    C4.5 & 1993 & \checkmark & \checkmark & \checkmark & \checkmark & \ding{55} \\
    TAO & 2018 & \checkmark & \checkmark & \ding{55} & \ding{55} & \ding{55} \\
    DL8.5 & 2020 & \checkmark & \checkmark & \checkmark & \ding{55} & \ding{55} \\
    GOSDT & 2022 & \checkmark & \checkmark & \ding{55} & \ding{55} & \ding{55} \\
    HS & 2022 & \checkmark & \checkmark & \ding{55} & \ding{55} & \ding{55} \\
    logicDT & 2024 & \ding{55} & \ding{55} & \checkmark & \ding{55} & \ding{55} \\
    \midrule
    \multicolumn{7}{l}{\textit{Rule List}} \\
    \midrule
    OneR & 1993 & \checkmark & \checkmark & \checkmark & \checkmark & \ding{55} \\
    BRL & 2015 & \ding{55} & \checkmark & \checkmark & \ding{55} & \ding{55} \\
    \midrule
    \multicolumn{7}{l}{\textit{Rule Set}} \\
    \midrule
    RIPPER & 1995 & \checkmark & \checkmark & \checkmark & \checkmark & \ding{55} \\
    Slipper & 1999 & \checkmark & \checkmark & \ding{55} & \ding{55} & \ding{55} \\
    M5Rules & 1999 & \checkmark & \ding{55} & \checkmark & \checkmark & \ding{55} \\
%    LRI & 2000 & \checkmark & \ding{55} & \ding{55} & ? & \ding{55} \\
    RuleFit & 2008 & \ding{55} & \checkmark & \ding{55} & \ding{55} & \ding{55} \\
    MLRules & 2008 & \checkmark & \ding{55} & \ding{55} & \checkmark & \ding{55} \\
    RuleKit & 2020 & \checkmark & \checkmark & \checkmark & \checkmark & \ding{55} \\
    \midrule
    \multicolumn{7}{l}{\textit{Evolutionary Computation}} \\
    \midrule
    GPAS & 2007 & \ding{55} & \ding{55} & \checkmark & \ding{55} & \checkmark \\
    BioHEL & 2009 & \checkmark & \ding{55} & \ding{55} & \ding{55} & \ding{55} \\
    ExSTraCS & 2015 & \checkmark & \checkmark & \ding{55} & \ding{55} & \ding{55} \\
    logicGP & 2025 & \checkmark & \ding{55} & \ding{55} & \ding{55} & \checkmark \\       
    \bottomrule
  \end{tabular}  
\end{table}

We can exclude all algorithms that do not directly model multiclass classification problems, as this is a key aspect of our analysis. We will also exclude algorithms that do not have open source implementations in popular machine learning libraries, as this would make a fair and transparent comparison difficult. Based on these criteria, the remaining candidates are: CART, C4.5, TAO, DL8.5, GOSDT, HS, OneR, RIPPER, Slipper, M5Rules, MLRules, RuleKit, ExSTraCS, and logicGP.

Many of the mentioned algorithms are available in imodels \citep{imodels2021}, which is a Python library for interpretable machine learning models. We try to choose only one algorithm where algorithms are expected to give similar results, e.g., CART and C4.5, to keep the comparison manageable. Based on this, we will compare the following algorithms in our experiments: CART, HS, OneR, Slipper, ExSTraCS, RuleKit, and logicGP. ExSTraCS \citep{urbanowicz_exstracs_2015}, RuleKit \citep{gudys2025rulekit2fastersimpler}, and logicGP \citep{10.1145/3712255.3734300} are not included in imodels, but they have open source implementations available, which allows us to include them in our comparison.

We will evaluate the performance of these algorithms on a set of benchmark UCI datasets. The evaluation will be based on F1 score, interpretability, and computational efficiency. The results of the comparison will be presented in the following sections.

\subsection{Experimental Setup}

The experiments are conducted on selected UCI benchmark datasets for classification. For each dataset, we apply the same preprocessing pipeline for all methods and choose the train-test split by dataset size: test size $0.30$ for $n<500$, $0.25$ for $500 \leq n < 5000$, and $0.20$ for $n \geq 5000$. To reduce variance due to random initialization and sampling effects, each experiment is repeated over 10 runs with different random seeds, and results are aggregated across runs. We report F1 score as the predictive metric together with interpretability-oriented model-size characteristics (rule, list, or tree size proxies) and discuss computational effort on a comparative level. We do not claim formal statistical significance beyond these repeated-run aggregates.

\subsection{UCI Datasets}

The UCI Machine Learning Repository \citep{uci_ml_repo} is a popular source of benchmark datasets for machine learning research. We select a small set of established classification benchmarks that cover binary and multiclass problems, different numbers of features, and different sample sizes. This gives a compact but diverse basis for comparing predictive performance and model complexity across the considered interpretable learning methods.

\begin{table*}
  \centering
  \caption{Overview of the UCI datasets used in the experiments}
  \label{tab:uci_datasets}
  \begin{tabular}{l l r r r}
    \toprule
    ID & Dataset & Features & Observations & Classes \\
    \midrule
    12 & Balance Scale & 4 & 625 & 3 \\
    267 & Banknote Authentication & 4 & 1372 & 2 \\
    17 & Breast Cancer Wisconsin Diagnostic & 30 & 569 & 2 \\
    19 & Car Evaluation & 6 & 1728 & 4 \\
    45 & Heart Disease & 13 & 303 & 5 \\
    53 & Iris & 4 & 150 & 3 \\
    78 & Page Blocks Classification & 10 & 5473 & 5 \\
    109 & Wine & 13 & 178 & 3 \\
    \bottomrule
  \end{tabular}
\end{table*}

Figure~\ref{fig:uci_comparison} shows the results of the comparison of the selected algorithms on the UCI datasets. The figure includes the F1 score and model size for each algorithm on each dataset. The results are aggregated across 10 runs with different random seeds to provide a more robust comparison. Note that Slipper sometimes completely fails to produce a model. Overall, the UCI benchmark results indicate meaningful performance and interpretability differences across the considered model families. In the next section, we interpret these patterns in more detail and derive implications for a unifying hypothesis-space perspective.

\begin{figure*}
  \centering
  \includegraphics[width=\textwidth]{merged_ucimodels_combined}  
  \caption{Comparison of algorithms on UCI datasets}
  \label{fig:uci_comparison}
\end{figure*}




\section{\uppercase{Discussion}}

The computational complexity is not given in the results, as it is not a key aspect of our analysis. However, it is worth noting that the algorithms that are based on evolutionary computation techniques, such as ExSTraCS and logicGP, are generally more computationally expensive than the other algorithms, as they require a large number of evaluations to find good solutions. However, the good performance of these algorithms often justifies the computational cost.

As visible in Figure~\ref{fig:uci_comparison}, Slipper and OneR can be largely excluded from the discussion of strong-performing candidates, as they perform worse than most of the other considered algorithms in our runs. With regard to interpretability, CART and ExSTraCS often produce comparatively large and complex models.

However, HS, which performs post-hoc optimization of decision trees, can often produce much smaller and therefore more interpretable models than CART while maintaining good performance. This suggests that strong predictive quality and high interpretability are not necessarily conflicting objectives. It also hints at the possibility of using better post-hoc optimization techniques for ExSTraCS, where existing rule-compaction procedures are not primarily tuned for interpretability.

From the perspective of this paper, these patterns are important because they become directly comparable within a single scored-rule-set view, rather than only as algorithm-specific outcomes.

Across our runs, RuleKit and logicGP tend to provide the best compromise of performance and interpretability in Figure~\ref{fig:uci_comparison}, often with small but accurate models. Both algorithms are rule-set-based: RuleKit follows a more traditional rule-set approach, while logicGP uses a broader scored-rule representation that can model more flexible decision boundaries. This supports the central claim that highly interpretable models can remain competitive on non-trivial multiclass problems.

Taken together, these empirical observations motivate a formal model-space view in which different interpretable learners can be expressed within one common representation. The next section makes this step explicit by introducing scored rule sets as a unifying hypothesis space.


\section{\uppercase{Scored Rule Sets}}

In the following, we will introduce a generalization of rule sets, which we call scored rule sets. This generalization allows us to address the potential disadvantages of rule sets while still maintaining their advantages. The idea of rule set scoring was introduced by \cite{10.1007/BFb0017011} and is e.g.\ used in logicGP \cite{10.1145/3712255.3734300}. We introduce a more general definition of scored rule sets than the one used in these works, as this allows us to gain more flexibility without losing interpretability. We will show in this section that scored rule sets 

\begin{itemize}
    \item generalize rule sets, rule lists, and decision trees, which allows them to represent a wider range of decision boundaries, and
    \item can address the potential disadvantages of unscored rule sets, such as contradictory rules and unclassified instances, while still maintaining interpretability.
\end{itemize}

%The problem of contradictory and missing rules can be addressed by using scored rule sets, where each rule is associated with a score that indicates the strength of the rule. The final prediction can then be made by combining the scores of all applicable rules, for example by taking a weighted average or by using a voting scheme. This approach can also enhance interpretability, as the scores can provide insights into the importance of each rule in the decision-making process.

In conclusion to the mentioned advantages and disadvantages, such a hypothesis space offers a good compromise of flexibility and interpretability. Formulated as a generalization of rule sets, it can represent a wide range of decision boundaries while still providing insights into the decision-making process through the associated scores.

\subsection{Definition of Scored Rule Sets}

We are given a database of observations described with a fixed number of attributes $\mathbf{A}_1, \ldots, \mathbf{A}_n$ and a designated class attribute $\mathbf{C} \in \{1, \ldots, m\}$. The learning task is to find a mapping $f$ such that
\begin{align*}
    f(\mathbf{A}_1, \ldots, \mathbf{A}_n) &\approx \bigl(P(\mathbf{C}=1 \mid \mathbf{A}_1,\ldots,\mathbf{A}_n), \\
    &\quad \ldots, P(\mathbf{C}=m \mid \mathbf{A}_1,\ldots,\mathbf{A}_n)\bigr)
\end{align*}
i.e.\ a function that estimates the probability of each class from the attribute values of new, previously unseen observations.

We define such a mapping via a \emph{scored rule set}. Note that our definition of a scored rule set is more general than the hypothesis spaces of e.g.\ \cite{10.1007/BFb0017011,10.1145/3712255.3734300}. This is motivated by gaining more flexibility without losing interpretability, as the more general hypothesis space can still be easily interpreted by humans.

Let $\mathcal{V}_j$ denote the domain of attribute $\mathbf{A}_j$. An \emph{atom} for attribute $\mathbf{A}_j$ is a function
\begin{equation*}
  a: \mathcal{V}_j \rightarrow \{\text{true}, \text{false}\}.
\end{equation*}
Let $\mathcal{A}$ denote the set of all possible atoms. A \emph{rule} $r = (S_r, \mathbf{s}_r)$ consists of a set of atoms $S_r \subseteq \mathcal{A}$ and a score vector $\mathbf{s}_r \in \mathbb{R}^m$. The pair induces a function
\begin{equation*}
  r: \mathcal{V}_1 \times \cdots \times \mathcal{V}_n \rightarrow \mathbb{R}^m,
\end{equation*}
where $r(\mathbf{x}) = \mathbf{s}_r$ if $a(x_j) = \text{true}$ for all $a \in S_r$ where $a$ is an atom for attribute $\mathbf{A}_j$, and $r(\mathbf{x}) = \mathbf{0}$ otherwise. Note that a rule with $S_r = \emptyset$ fires for every observation, since the condition is vacuously true. Such a \emph{default rule} can be included in $R$ to ensure a well-defined prediction even when no other rule applies; its score vector $\mathbf{s}_{r_0}$ can encode a prior over the classes, e.g., estimated from class frequencies in the training data.
Let $\mathcal{R}$ denote the set of all rules. A \emph{scored rule set} is a finite set $R \subseteq \mathcal{R}$ together with an \emph{aggregation function}
\begin{equation*}
  g: 2^{\mathcal{R}} \times (\mathcal{V}_1 \times \cdots \times \mathcal{V}_n) \rightarrow [0,1]^m,
\end{equation*}
which maps a rule set and an observation to a vector of class probabilities. A common choice is
\begin{equation*}
  g(R, \mathbf{x})[i] = \frac{\exp\!\left(\sum_{r \in R} r(\mathbf{x})[i]\right)}{\sum_{k=1}^{m} \exp\!\left(\sum_{r \in R} r(\mathbf{x})[k]\right)},
\end{equation*}
i.e.\ the softmax of the cumulative scores across all rules in $R$. Another common choice is
\begin{equation*}
  g(R, \mathbf{x})[i] = \begin{cases} 1 & \text{if } i = \arg\max_{k \in \{1,\ldots,m\}} \sum_{r \in R} r(\mathbf{x})[k], \\ 0 & \text{otherwise,} \end{cases}
\end{equation*}
i.e.\ assigning full probability to the class with the highest cumulative score.


\subsection{Relation to Other Models}

Scored rule sets can be seen as a generalization of rule sets in which each rule contributes a class-score vector. This supports a flexible multiclass representation and allows the hypothesis spaces of the algorithms discussed above to be modeled as special cases. Outputs that are not score vectors can be one-hot encoded or mapped to scores, and the aggregation function can mirror the original prediction mechanism (e.g., voting, $\arg\max$, or softmax).

%\subsubsection{Relation to Rule Lists and Decision Trees}

As noted above, rule sets can represent any decision tree or rule list (see e.g.\ \cite{witten2025data}); therefore scored rule sets inherit this property. In addition, score magnitudes can encode structural information such as rule priority in rule lists or depth-based emphasis in decision trees.

% \subsubsection{Relation to other Rule Sets}

% The discussed scored rule sets can also be seen as a generalization of the rule sets used in RuleKit and the rule populations used in Learning Classifier Systems (LCS), which are a type of evolutionary algorithm that learns a population of rules to solve classification problems. In LCS, each rule is associated with a fitness value that indicates the strength of the rule, and the final prediction is made by combining the predictions of all applicable rules based on their fitness values. This is similar to the scored rule sets, where each rule is associated with a score vector that indicates the strength of the rule for each class, and the final prediction is made by combining the scores of all applicable rules.

\subsection{Examples}

Our examples are oriented towards the most successful algorithms in our experiments: HS, ExSTraCS, RuleKit, and logicGP. 

The hypothesis space of scored rule sets is inspired by the one used in logicGP \citep{10.1145/3712255.3734300}, which is a genetic programming algorithm for learning scored rule sets for multiclass classification problems. logicGP models can be transformed almost directly into the more general scored rule sets defined above and are therefore a good example to illustrate the concept of scored rule sets.

\subsubsection{logicGP / RuleKit}

Table~\ref{tab:iris_example} shows an exemplary scored rule set $(R, g)$ for the Iris dataset with $m=3$ classes: setosa ($i=1$), versicolor ($i=2$), and virginica ($i=3$). Each row represents a rule $r = (S_r, \mathbf{s}_r) \in R$, where $S_r$ is the set of atoms and $\mathbf{s}_r[1], \mathbf{s}_r[2], \mathbf{s}_r[3]$ are the corresponding class scores. The first rule is the default rule with $S_r = \emptyset$.

\begin{table}
   \centering
  \caption{Exemplary logicGP model for the Iris dataset.}
   \label{tab:iris_example}
   \begin{tabular}{p{0.5cm} p{0.5cm} p{0.5cm} p{4.2cm}}
      \toprule
      $\mathbf{s}_r[1]$ & $\mathbf{s}_r[2]$ & $\mathbf{s}_r[3]$ & $S_r$ \\
      \midrule
      $1.00$ & $0.00$ & $0.00$ & $\emptyset$ \\
      $0.00$ & $2.31$ & $1.39$ & $\{\text{petal width} \in [1.4,\,1.6]\}$ \\
      $0.00$ & $5.56$ & $0.00$ & $\{\text{sepal width} \in [2.0,\,3.0],\ $ $\text{petal width} \in [0.4,\,1.3]\}$ \\
      $0.00$ & $0.00$ & $7.00$ & $\{\text{petal length} \in [5.2,\,6.9]\}$ \\
      $0.00$ & $0.00$ & $7.00$ & $\{\text{petal width} \in [1.9,\,2.5]\}$ \\
      \bottomrule
   \end{tabular}
\end{table}

%\citep{Unwin2021TheID}

The default rule ($S_r = \emptyset$, first row) fires for every observation and acts as a prior; the remaining rules each define a non-empty set of atoms $S_r$ whose conjunction must hold for the score $\mathbf{s}_r$ to be added to the cumulative class scores.

RuleKit's rule lists can be easily transformed into scored rule sets by assigning appropriate scores to the rules and using an aggregation function that reflects the way predictions are made in RuleKit, which is a voting scheme. 

\subsubsection{ExSTraCS} 

\citet{urbanowicz_exstracs_2015} introduce ExSTraCS, an evolutionary algorithm that learns scored rule sets for classification problems. ExSTraCS uses a genetic algorithm to evolve a population of rules. An excerpt of rules for an exemplary run on the Iris dataset are shown in Table~\ref{tab:exstracs_iris_example}. Each row represents a rule with conditions on the attributes ($s_l=\text{sepal length}$, $s_w=\text{sepal width}$, $p_l=\text{petal length}$, $p_w=\text{petal width}$) and a corresponding class label ($0$ for setosa, $1$ for versicolor, and $2$ for virginica) and a score ($f\times n = \text{fitness} \times \text{numerosity}$) that indicates the strength of the rule. An $\#$ indicates that the corresponding attribute is not used in the rule, which means that it always evaluates to true for that attribute. The first three rules are default rules that fire for every observation and act as priors for the corresponding class. The remaining rules define conditions on the attributes that must hold for the rule to fire and contribute to the final prediction.

\begin{table}
   \centering
  \caption{Excerpt from ExSTraCS model for the Iris dataset.}
   \label{tab:exstracs_iris_example}
   \begin{tabular}{p{0.6cm} p{0.6cm} p{0.6cm} p{1.3cm} p{0.8cm} p{0.8cm}}
      \toprule
      $s_l$ & $s_w$ & $p_l$ & $p_w$ & class & $f\times n$ \\
      \midrule
    \# & \# & \# & \# & $2$ & $13$ \\
    \# & \# & \# & \# & $0$ & $10$ \\
    \# & \# & \# & \# & $1$ & $9$ \\
    \# & \# & \# & $0.74,\ 2.10$ & $1$ & $6$ \\
    %\# & \# & \# & $-0.14,0.58$ & $0$ & $6$ \\
    %\# & \# & \# & $0.02,0.51$ & $0$ & $5$ \\
    %\# & \# & \# & $-0.14,0.51$ & $0$ & $5$ \\
    %\# & \# & $0.41,2.13$ & \# & $0$ & $5$ \\
    %$4.50,7.10$ & \# & \# & \# & $2$ & $5$ \\
    %\# & \# & $0.41,2.42$ & \# & $0$ & $5$ \\
    %\# & \# & $3.92,7.45$ & \# & $2$ & $5$ \\
    %\# & \# & $0.41,2.06$ & \# & $0$ & $5$ \\
    %\# & \# & $4.37,7.46$ & \# & $2$ & $5$ \\
    %\# & \# & \# & $-0.14,0.73$ & $0$ & $5$ \\
      \bottomrule
   \end{tabular}
\end{table}

ExSTraCS uses an alternative scored-rule-set representation in which rule strength is encoded through fitness and numerosity, while conditions are expressed as intervals. Mapping this representation to our definition is straightforward: intervals correspond to atoms, and fitness-numerosity information can be interpreted as class scores. Using score vectors also allows a more compact representation than standard ExSTraCS, where each rule is tied to only one class label and a single score.
  
Thus, a corresponding scored rule set is given in Table~\ref{tab:exstracs_iris_example_transformed}, where each row represents a rule $r = (S_r, \mathbf{s}_r) \in R$, with $S_r$ being the set of atoms and $\mathbf{s}_r[1], \mathbf{s}_r[2], \mathbf{s}_r[3]$ being the corresponding class scores.

\begin{table}
   \centering
  \caption{Transformed ExSTraCS model for the Iris dataset.}
   \label{tab:exstracs_iris_example_transformed}
   \begin{tabular}{p{0.6cm} p{0.6cm} p{0.6cm} p{4.2cm}}
      \toprule
      $\mathbf{s}_r[1]$ & $\mathbf{s}_r[2]$ & $\mathbf{s}_r[3]$ & $S_r$ \\
      \midrule
      $10.00$ & $9.00$ & $13.00$ & $\emptyset$ \\
      $0.00$ & $6.00$ & $0.00$ & $\{\text{petal width} \in [0.74,\,2.10]\}$ \\
      %$16.00$ & $0.00$ & $0.00$ & $\{\text{petal width} \in [-0.14,\,0.58]\}$ \\
      %$10.00$ & $0.00$ & $0.00$ & $\{\text{petal length} \in [0.41,\,2.42]\}$ \\
      %$0.00$ & $0.00$ & $15.00$ & $\{\text{petal length} \in [3.92,\,7.46]\}$ \\
      %$5.00$ & $0.00$ & $0.00$ & $\{\text{petal length} \in [0.41,\,2.06]\}$ \\
      %$6.00$ & $0.00$ & $0.00$ & $\{\text{petal width} \in [-0.14,\,0.73]\}$ \\
      \bottomrule
   \end{tabular}
\end{table}

  
  %In the example, the first three rules could therefore be merged into a single rule.

\subsubsection{CART / HS}

CART \citep{Breiman1984} constructs a decision tree by recursively selecting splits that improve class separation, while HS \citep{pmlr-v162-agarwal22b} provides post-hoc optimization for more compact and interpretable trees. Figure~\ref{fig:cartdt} shows an Iris example; leaves are displayed with one-hot class-probability vectors, consistent with an $\arg\max$ decision rule.

\begin{figure}
  \centering
  \includegraphics[width=\columnwidth]{decisiontree_iris_compact}  
  \caption{Example of a CART decision tree for the Iris dataset.}
  \label{fig:cartdt}
\end{figure}

The tree can be transformed into a scored rule set by mapping each root-to-leaf path to a rule with class-score vector. This enables direct comparison with logicGP, RuleKit, and ExSTraCS models in one representation. Depth-dependent scoring can encode structural priority and support atom removal while preserving the predicted class under $\arg\max$.

% This approach is discussed in more detail in the following section, where we will also introduce an exact aggressive pruning procedure to further simplify the resulting scored rule set while preserving the final class prediction under an $\arg\max$ decision rule.

% \subsubsection{Scored Rule Set Extraction from Decision Trees}

% From this construction, we derive a compact scored rule set directly from the tree while preserving the final class prediction under an $\arg\max$ decision rule.

% \paragraph{Pruning-focused weighted rule extraction.}
% We focus on the practically important case that the input tree has no non-leaf class-homogeneous subtree. In this setting, subtree-level compression cannot reduce the rule count, so the main simplification mechanism is atom pruning. The unified procedure combines (i) depth-weighted leaf-to-rule conversion and (ii) exact aggressive atom pruning with full equivalence validation on the tree-induced partition.

% Let $T$ be a multiclass decision tree with class set $\{1,\ldots,m\}$ and leaf set $L(T)$. For each leaf $\ell\in L(T)$ at depth $d(\ell)$ with class $c(\ell)$, create one rule
% \begin{equation*}
%   r_\ell=(S_\ell,\mathbf{s}_\ell),
% \end{equation*}
% where $S_\ell$ is the conjunction of split decisions on the path root$\to\ell$, and
% \begin{equation*}
%   \mathbf{s}_\ell[i] = \begin{cases}
%     \lambda^{-d(\ell)} & \text{if } i=c(\ell),\\
%     0 & \text{otherwise,}
%   \end{cases}
%   \qquad \lambda>1.
% \end{equation*}
% Then, iterate over atoms in the rules and accept a removal only if prediction equivalence is preserved on the complete leaf-region partition of the original tree.

% \begin{algorithm}[t]
% \caption{Exact aggressive pruning-focused depth-weighted conversion}
% \label{alg:depth_weighted_tree_compression}
% \KwIn{Multiclass decision tree $T$, decay parameter $\lambda>1$}
% \KwOut{Depth-weighted and atom-pruned scored rule set $R$}

% $R \leftarrow \emptyset$
% \;
% \ForEach{leaf $\ell \in L(T)$}{
%   $S_\ell \leftarrow$ conjunction of split tests on path root$\to\ell$
%   \;
%   initialize $\mathbf{s}_\ell$ with one non-zero entry: $\mathbf{s}_\ell[c(\ell)] \leftarrow \lambda^{-d(\ell)}$
%   \;
%   $R \leftarrow R \cup \{(S_\ell,\mathbf{s}_\ell)\}$
%   \;
% }

% $\mathcal{L} \leftarrow$ leaf-region partition induced by $T$
% \;
% $\texttt{changed} \leftarrow \texttt{true}$
% \;
% \While{$\texttt{changed}$}{
%   $\texttt{changed} \leftarrow \texttt{false}$
%   \;
%   \ForEach{rule $r=(S_r,\mathbf{s}_r) \in R$}{
%     \ForEach{atom $a \in S_r$}{
%       build candidate $R'$ by replacing $r$ with $r'=(S_r\setminus\{a\},\mathbf{s}'_r)$
%       \;
%       \If{$\arg\max$ prediction is unchanged on all regions in $\mathcal{L}$}{
%         $R \leftarrow R'$
%         \;
%         $\texttt{changed} \leftarrow \texttt{true}$
%         \;
%       }
%     }
%   }
% }
% \Return $R$
% \end{algorithm}

% The initialization phase is linear in the tree size, i.e., $\mathcal{O}(|V|)$ for node set $V$. The pruning phase dominates the runtime and is data-independent in the sense that validation is performed on the finite region partition $\mathcal{L}$ induced by the tree. A practical bound is
% \begin{equation*}
%   \mathcal{O}(I\cdot |R|\cdot \bar{k}\cdot |\mathcal{L}|\cdot m),
% \end{equation*}
% where $I$ is the number of outer iterations, $\bar{k}$ the average number of atoms per rule, and $m$ the number of classes.

% Prediction preservation follows directly from the acceptance criterion: each accepted pruning step keeps the same $\arg\max$ class on every region in $\mathcal{L}$, hence by induction over accepted updates the final rule set remains equivalent to the original tree on all inputs.

% For the concrete Iris tree in Figure \ref{fig:cartdt}, this focus is particularly appropriate: the tree contains no compressible class-homogeneous non-leaf subtree, so rule-count reduction is achieved by atom pruning rather than subtree merging. The initial weighted model still contains $9$ rules (one per leaf), with weights $\lambda^{-1}$, $\lambda^{-3}$, $\lambda^{-4}$, and $\lambda^{-5}$ depending on leaf depth.

% For example, choosing $\lambda=2$ yields the concrete depth weights $2^{-1}=0.50$, $2^{-3}=0.125$, $2^{-4}=0.0625$, and $2^{-5}=0.03125$. This keeps the same predicted class under $\arg\max$ while making shallow rules numerically more influential.

Applying such a method to the Iris tree (with full validation on all tree-induced regions) yields the equivalent compact representation in Table~\ref{tab:cart_iris_example_depth_weighted_pruned} when $\lambda=2$. In this run, $22$ atoms are removed while preserving exact prediction equivalence. Note that this compaction is only possible because of the depth-weighted scoring, which allows the pruning procedure to remove atoms while keeping the same predicted class under $\arg\max$.

%The number of removable atoms depends on the depth-weighting parameter $\lambda$; for the Iris tree considered here, larger $\lambda$ values allowed more removals up to a plateau. Whether this behavior holds more generally is beyond the scope of the present section. A global maximization of removable atoms can be formulated as an optimization problem, but this is likewise not our focus here.



\begin{table}
  \centering
  \caption{Depth-weighted scored rule set for Iris}
  \label{tab:cart_iris_example_depth_weighted_pruned}
  \begin{tabular}{p{0.6cm} p{0.6cm} p{0.6cm} p{3.2cm}}
    \toprule
    $\mathbf{s}_r[1]$ & $\mathbf{s}_r[2]$ & $\mathbf{s}_r[3]$ & $S_r$ \\
    \midrule
    $\lambda^{-1}$ & $0$ & $0$ & $\{\text{petal length} \leq 2.45\}$ \\
    $0$ & $\lambda^{-4}$ & $0$ & $\{\text{petal length} \leq 4.95,\ $ $\text{petal width} \leq 1.65\}$ \\
    $0$ & $0$ & $\lambda^{-4}$ & $\{\text{petal length} \leq 4.95\}$ \\
    $0$ & $0$ & $\lambda^{-4}$ & $\{\text{petal length} > 4.95,\ $ $\text{petal width} \leq 1.55\}$ \\
    $0$ & $\lambda^{-5}$ & $0$ & $\emptyset$ \\
    $\lambda^{-5}$ & $0$ & $0$ & $\{\text{petal length} > 5.45\}$ \\
    $0$ & $\lambda^{-4}$ & $0$ & $\{\text{petal width} > 1.75\}$ \\
    $0$ & $0$ & $\lambda^{-4}$ & $\{\text{petal width} > 1.75,\ $ $\text{sepal length} > 5.95\}$ \\
    $0$ & $0$ & $\lambda^{-3}$ & $\{\text{petal width} > 1.75,\ $ $\text{petal length} > 4.85\}$ \\
    \bottomrule
  \end{tabular}
\end{table}

%No score re-optimization is applied in this variant; accepted updates are based solely on exact prediction preservation on $\mathcal{L}$.

\section{\uppercase{Future Work}}

Several directions appear particularly promising for future work. First, better rule-compaction procedures for ExSTraCS could improve interpretability without sacrificing predictive quality. Second, integrating logicGP into the scikit-learn ecosystem would substantially improve accessibility, reproducibility, and comparability in practical workflows. Third, although logicGP already learns in a closely related scored-rule-set space, it would be interesting to investigate whether new learning procedures in this unified space can perform even better by integrating search strategies, optimization principles, or compaction techniques from other research communities. Finally, the scored-rule-set perspective suggests the development of a dedicated scikit-learn-compatible package that exposes scored rule sets as a common target representation and includes transformations from multiple underlying algorithms into this unified hypothesis space.

\section{\uppercase{Conclusion}}

This work contributes a unifying perspective on interpretable multiclass classification by connecting model classes that are often treated separately in the literature. We showed that scored rule sets provide a common and interpretable representation that subsumes classical rule sets and can also encode rule lists and decision-tree structures.

Empirically, our benchmark comparison on selected UCI datasets indicates that strong predictive performance and compact model representations can be achieved within interpretable model families, with rule-set-based methods often providing a favorable trade-off. Conceptually, the citation and methodology analysis suggests that limited cross-community integration still exists, particularly between evolutionary and non-evolutionary research lines, and that a shared hypothesis-space view can help reduce this fragmentation.

The paper also provided constructive model-level mappings from logicGP, ExSTraCS, and CART representations into scored rule sets, illustrating that the framework is not only theoretical but directly applicable to existing methods. These examples highlight how scoring mechanisms can resolve typical rule-set issues such as conflicting evidence and uncovered instances while preserving transparency.

%A limitation of the present study is the finite set of datasets and algorithms, and the absence of a dedicated runtime-complexity benchmark in the empirical section.

\bibliographystyle{ACM-Reference-Format}
\bibliography{scoredrulesets} 


\end{document} 